\activity{9}{Observability}{observability}{Saturday, 10:04, four minutes into the stadium on-sale. Support fills with ``checkout spins forever''. Every pod is green, CPU sits at 45\%, the error rate is 0.3\%, the health checks pass, and yet the median fan waits nine seconds to pay. Four teams each show a dashboard proving that \emph{their} service is fine. Which signals, shared context and alerts would have found the slow hop in minutes, and who acts on them?}
\stages{\stClarify\ \stEstimate\ \stDeep}

\begin{know}
  \item \textbf{Monitoring} answers questions you asked in advance; \textbf{observability} lets you infer the internal state from the outputs and ask new questions. Five palette components, one question each: metrics (is something wrong?), logs (what happened?), tracing (where?), alerting (must a human act now?), health check (should this instance get traffic?).
  \item \textbf{RED} (rate, errors, duration) for services; \textbf{USE} (utilization, saturation, errors) for resources such as CPU, pools and queues. Saturation (work waiting) is the leading indicator.
  \item \textbf{Cardinality}: the number of series of a metric is the product of the number of values of each label. Keep labels bounded; per-customer detail goes to traces and logs.
  \item Latency goes in \textbf{histograms}: sum the buckets of all pods (\texttt{sum by (le)}) before computing a percentile. Never average percentiles.
  \item A \textbf{trace} is a tree of spans under one 128-bit \texttt{trace\_id}, propagated with the W3C \texttt{traceparent} header; logs carry the same id. Spans for \emph{waiting} (pool, lock, queue) reveal saturation.
  \item \textbf{Sampling}: head sampling decides blindly at the start; tail sampling (in the OpenTelemetry Collector) keeps every error and slow trace. Metrics are computed from 100\% of the traffic, before sampling.
  \item \textbf{Burn rate} $b = \text{error ratio} / (1 - \text{SLO})$; budget consumed $= b \cdot t / \text{window}$. Multi-window rules: page at $14.4\times$ over 1 h (confirmed by 5 min) and $6\times$ over 6 h (30 min); ticket at $1\times$ over 3 days (6 h). Alert on symptoms, not causes.
  \item \textbf{Incidents}: mitigate first, then investigate. MTTD (detect), MTTM (mitigate) and MTTR (resolve) count from the start of the impact, MTTA (acknowledge) from the page; blameless postmortems end in owned, dated action items.
\end{know}

\begin{ex}{1}{Which signal answers?}
\items
  \item Name the palette component that answers each question: (1) ``the checkout error rate jumped from 0.1\% to 4\%: is something wrong, and how much?''; (2) ``3.8 s of a 4 s checkout: where?''; (3) ``what exactly happened to order 81234?''; (4) ``the checkout SLO is burning 14$\times$ faster than sustainable: must someone act now?''; (5) ``\texttt{/health/ready} returns 503 during warm-up: should this pod get traffic?''.
  \item Classify each metric as RED or USE: the request rate of \texttt{POST /checkout}; the CPU utilization of the inventory database; the pending acquisitions of the payment connection pool; the 5xx ratio of \texttt{POST /checkout}; the p99 of \texttt{POST /checkout}. Which one is a saturation signal?
  \item In the hook, CPU at 45\% and errors at 0.3\% looked fine. Why are these the wrong signals to watch for ``checkout spins forever''?
\end{enumerate}
\lines{5}
\end{ex}

\simulator{You can't fix what you don't measure}{observability-stack}{Utilization shows ``?'' until a service is connected to Metrics. Wire Metrics, a Health Check on the load balancer and Alerting, then compare how long the replica crash at t = 50 s costs with and without them.}

\begin{ex}{1}{Counting series}
A latency histogram has 15 buckets plus \texttt{\_sum} and \texttt{\_count}. Each series costs about 4 KiB of memory in the metrics store. Each row adds one label to the previous ones.
\begin{center}\small
\begin{tabular}{@{}lrrr@{}}
  \toprule
  label added & values & series & memory \\
  \midrule
  histogram (15 buckets $+$ \texttt{\_sum} $+$ \texttt{\_count}) & 17 & \blank\hspace{1.8cm} & \hspace{1.8cm} \\
  \texttt{http\_request\_method} & 4 & \blank & \\
  \texttt{http\_route} & 25 & \blank & \\
  \texttt{http\_response\_status\_code} & 6 & \blank & \\
  \texttt{pod} & 20 & \blank & \\
  \texttt{customer\_id} & 1\,000\,000 & \blank & \\
  \bottomrule
\end{tabular}
\end{center}
\items
  \item Fill in the table.
  \item A developer wants \texttt{customer\_id} as a label ``to debug angry fans''. Where does that information belong instead?
\end{enumerate}
\lines{3}
\end{ex}

\begin{ex}{2}{Four pods, one p99}
The fleet behind \texttt{POST /checkout} served 100\,000 requests in an hour. Pods \texttt{a}, \texttt{b} and \texttt{c} served 30\,000 each with a p99 of 180 ms. Pod \texttt{d}, on a noisy host, served the other 10\,000 with a median of 400 ms, a p90 of 760 ms and a p99 of 1\,260 ms.
\items
  \item A dashboard averages the four p99s. What does it show?
  \item How many requests are in the slowest 1\% of the fleet? Where do they come from? Estimate the true fleet p99.
  \item The histogram buckets are 0.1, 0.25, 0.5, 1 and 2.5 s (cumulative, ``less than or equal''). Which buckets does a 0.3 s request increment? The SLO is ``checkouts under 2 s'': can you compute it from these buckets?
  \item Write, in words, how the fleet p99 must be computed.
\end{enumerate}
\lines{6}
\end{ex}

\begin{ex}{2}{What telemetry costs}
\items
  \item A normal day: 500 req/s, each producing 8 spans of 500 B and 5 log lines of 300 B. Compute the trace and log volumes in MB/s and GB/day. And per 30 days of traces at 100\% sampling?
  \item A busy day: 2\,000 req/s, 8 spans of 500 B per request, 1\% of the requests fail and 2\% are slower than 2 s. Compute the trace storage per day and the error traces kept per day for: keep everything; head sampling at 10\%; tail sampling that keeps all errors, all slow requests and 5\% of the rest.
  \item The flash sale multiplies traffic by 50. Which of the three strategies would you run, and why must the RED metrics still be computed from all requests?
\end{enumerate}
\lines{6}
\end{ex}

\begin{ex}{2}{Burn rates and detection time}
The checkout SLO is 99.9\% over a rolling \textbf{28-day} window (672 h).
\items
  \item How many minutes of full outage does the budget allow?
  \item For an error ratio of 1.44\%: burn rate, share of the budget spent in 1 h, and time until the budget is gone. Same for 0.6\% (share spent in 6 h).
  \item The fast page fires when the 1-hour error ratio exceeds $14.4 \times 0.1\% = 1.44\%$. For a constant error ratio $e$ starting at time 0, it fires after $t = 60 \times 1.44\% / e$ minutes. Compute $t$ for $e = 100\%$ and $e = 2\%$. For $e = 0.7\%$, which rule fires (the 6 h rule uses $6 \times 0.1\% = 0.6\%$), and after how long?
  \item At 10:03 of an on-sale the error ratio jumps to 4\% and stays there. When does the fast page fire? The on-sale is over in 30 minutes: what do you change?
\end{enumerate}
\lines{6}
\end{ex}

\begin{ex}{2}{Read the waterfall}
The trace of one slow checkout (total 9.12 s): gateway 9.12 s $\to$ \texttt{order} 9.08 s $\to$ \texttt{inventory} \texttt{PUT /holds} 0.09 s (of which PostgreSQL 0.02 s); \texttt{payment} \texttt{POST /charges} 8.95 s, of which \texttt{pool.acquire} 7.90 s and the provider call 1.02 s; Kafka publish 0.01 s. The payment adapter has a pool of 20 connections.
\items
  \item Where is the critical path? What share of the 9.12 s is spent waiting inside BoxOffice?
  \item With Little's Law, how many charges per second can a pool of 20 connections sustain when each provider call takes 1.02 s? What happens to fans when more arrive?
  \item Fill in symptom $\to$ cause $\to$ fix: why did CPU at 45\% and green health checks hide the problem, and which two changes would you make in the code and which metric would you watch?
\end{enumerate}
\lines{6}
\end{ex}

\simulator{A slow dependency saturates a pool}{cascading-failure}{Run as is and read the bottleneck in the Report: threads exhausted while the CPU is almost idle. Then add a 300 ms timeout on the call to the external API and compare the error rate and p99 of the routes that do not use it.}

\begin{ex}{3}{The incident, the bill and the wiring}
The on-sale timeline: 10:00 sales open; 10:02 impact starts; 10:03 page; 10:04 acknowledged; 10:06 SEV1 declared; 10:11 an exemplar trace shows the pool wait; 10:14 mitigated; 10:40 resolved. About 12\% of the checkouts failed between 10:02 and 10:14, at 200 checkouts/s. In a normal 30-day window BoxOffice has about 2 million checkouts, with an SLO of 99.9\%.
\items
  \item Compute MTTD, MTTA, MTTM and MTTR for this incident.
  \item How many checkouts failed? How many months' worth of error budget is that?
  \item Draw the observability wiring of the checkout path: which component observes which, where the \texttt{trace\_id} flows, which alert pages and which opens a ticket.
  \item Write two postmortem action items (what, type, owner) that would have prevented or shortened this incident.
\end{enumerate}
\palette{\comp{edge}{API Gateway}\ \comp{compute}{App Server}\ \comp{storage}{SQL Database}\ \comp{external}{Payment}\ \comp{observability}{Metrics}\ \comp{observability}{Logs}\ \comp{observability}{Tracing}\ \comp{observability}{Alerting}\ \comp{observability}{Health Check}}
\frameblank{5.2cm}
\lines{3}
\end{ex}
